\section{推理运行时与内存规划：typed graph 如何服务一次真实请求}

\subsection{先从一个会抖动的请求开始}

模型导入和 verifier 把外部资产固定成 typed graph；AI 编译器把 typed graph 降成 executable plan。现在看计划真正落到一次请求时会遇到什么问题。假设用户输入 512 个 token 的 prompt，希望继续生成 128 个 token，本地引擎已经通过 verifier，权重也能被 kernel 读取。一个最朴素的 runtime 可能这样写：

\begin{lstlisting}[numbers=none,caption={朴素 runtime 的热路径}]
tokens = tokenizer.encode(prompt)

for op in graph:
  input_tensors = lookup_inputs(op)
  output_tensor = allocate(op.output_shape, op.output_dtype)
  run_kernel(op, input_tensors, output_tensor)

for step in 0..max_new_tokens:
  for op in graph:
    input_tensors = lookup_inputs(op)
    output_tensor = allocate(op.output_shape, op.output_dtype)
    run_kernel(op, input_tensors, output_tensor)
  next_token = sample(logits)
\end{lstlisting}

这段代码很容易在小模型上“跑通”，也很容易在 7B 级模型上暴露几个实际问题。第一，每个算子动态分配输出，会把 allocator、缺页、清零和碎片化放进热路径；decode 每生成一个 token 都重复这些成本，尾延迟会抖。第二，每次遍历 graph 时还在解释 op、查 dtype、查 layout、判断 shape，分支和间接调用压进最热循环。第三，所有 activation 都当作独立对象存活，内存峰值会被中间张量堆起来。第四，KV cache 是跨 token 状态，却被朴素代码混在普通临时张量里，容易出现覆盖、重复分配或错误增长。

真实推理运行时要做的不是“按图执行一遍”，而是把第 6、7 章的 typed graph、shape 合同、layout 合同和后端计划绑定到一个 session 状态上。编译期已经能决定的事情，不能留到每个 token 再判断；每个 token 才知道的事情，也不能假装编译期已经固定。

\begin{keyidea}
推理运行时的核心任务，是把编译器产出的 executable plan 绑定到一次会增长、会采样、会持有 KV cache 的真实请求。内存规划不是运行时小优化，而是让计划稳定执行的前提。
\end{keyidea}

\subsection{runtime state：请求不是一串独立 kernel}

一次本地对话请求至少包含下面这些运行期状态：

\begin{itemize}
  \item session：模型实例、后端实例、编译计划、线程池、设备句柄和错误报告上下文。
  \item token buffer：prompt token、已生成 token、当前位置、最大上下文和 stop 条件。
  \item KV cache：每层已经写入的 K/V、当前可读长度、容量、布局和可能的 page table。
  \item activation arena：prefill 或 decode 当前阶段可复用的临时张量区域。
  \item workspace：某些后端 kernel 需要的临时区，例如 attention scratch、GEMM packing scratch 或 GPU workspace。
  \item sampler state：temperature、top-k、top-p、seed、重复惩罚和 logits 后处理状态。
  \item profiler/oracle：记录事件、内存峰值、kernel 耗时和可选 reference 对比。
\end{itemize}

这些状态不属于同一种生命周期。权重从模型加载活到模型卸载；packed weight 从后端准备完成活到后端销毁；activation arena 通常只活在一个 plan invocation 内；KV cache 活在整个 session 内，并且随 token 增长；sampler state 随生成循环推进；profiler 事件则是运行期证据。把它们全部塞进一个“tensor map”会让内存所有权模糊，后续优化也没有安全边界。

\begin{lstlisting}[numbers=none,caption={推理 runtime 的核心对象关系}]
Engine
  ModelManifest
  TypedGraph
  ExecutablePlans
  BackendContext

Session
  TokenBuffer
  KVCache
  ActivationArena
  WorkspacePool
  SamplerState
  ProfilerTrace
\end{lstlisting}

注意这里的分层。Engine 持有模型级事实和编译计划，多个 session 可以共享它；Session 持有请求级状态，不能被另一个请求随意复用。CPU 后端的 packed weight、线程池和 NUMA 放置属于 Engine/BackendContext；某个用户这次对话的 KV cache 和 token buffer 属于 Session。这个边界一旦错了，多请求并发时会出现更隐蔽的问题：一个 session 覆盖另一个 session 的 cache，或者一个用户的 profiler/oracle 开关拖慢所有用户的热路径。

\subsection{内存分类：先分清谁能复用}

内存 planner 的第一步不是写分配器，而是给所有数据分类。下面这张表比“都叫 tensor”更接近真实运行时：

\begin{center}
\begin{tabular}{p{0.18\linewidth}p{0.23\linewidth}p{0.26\linewidth}p{0.22\linewidth}}
\toprule
类别 & 生命周期 & 典型位置 & 复用规则 \\
\midrule
原始权重 & 模型加载到卸载 & mmap 文件或主内存 & 只读，不能当 scratch \\
packed weight & 后端准备到卸载 & CPU 对齐内存或 GPU 显存 & 后端私有布局，按 kernel 合同读取 \\
activation & 一段执行计划内部 & arena 或 device buffer & last use 后可复用 \\
workspace & 单个 kernel 或 segment 内部 & workspace pool & 按算法临时借用，调用后归还 \\
KV cache & session 生命周期 & 主内存、显存或分层缓存 & 追加写入，历史可读，不能被普通 arena 覆盖 \\
logits & 每步采样前后 & 小 buffer 或 device-to-host staging & 可能需要保留 top-k 或完整分布 \\
\bottomrule
\end{tabular}
\end{center}

这张表背后是编译器理论里的 liveness analysis。一个 activation 的定义点是某个算子的输出，最后一次使用点是最后一个消费者；最后一次使用之后，它占用的内存槽可以给别的 activation。原始权重和 KV cache 不满足这个条件：权重会被每层、每步反复读取；KV cache 的旧 position 虽然不会再被写入，却会在后续 attention 中反复读取。因此它们不能进入普通 activation arena。

\begin{systemdiagram}{推理内存的生命周期分类}
\begin{pdfdiagram}[x=1cm,y=1cm]
  \node[book row label] at (-0.35,3.15) {模型级};
  \node[book accent node,text width=2.25cm] (w0) at (1.25,3.15) {原始权重\\mmap/load};
  \node[book teal node,text width=2.25cm] (w1) at (4.05,3.15) {packed\\weight};
  \node[book label,text width=2.4cm] at (7.05,3.15) {只读或后端私有\\跨请求共享};

  \node[book row label] at (-0.35,1.75) {请求级};
  \node[book purple node,text width=2.25cm] (kv) at (1.25,1.75) {KV cache\\跨 token};
  \node[book amber node,text width=2.25cm] (tok) at (4.05,1.75) {token/sampler\\状态};
  \node[book label,text width=2.4cm] at (7.05,1.75) {随 session 增长\\不能互相覆盖};

  \node[book row label] at (-0.35,0.35) {计划级};
  \node[book navy node,text width=2.25cm] (arena) at (1.25,0.35) {activation\\arena};
  \node[book red node,text width=2.25cm] (ws) at (4.05,0.35) {workspace\\scratch};
  \node[book label,text width=2.4cm] at (7.05,0.35) {按 last use\\和 kernel 调用复用};

  \draw[book strong arrow] (w0) -- (w1);
  \draw[book weak arrow] (kv) -- (tok);
  \draw[book weak arrow] (arena) -- (ws);

  \draw[book boundary] (-0.05,2.55) rectangle (5.45,3.75);
  \draw[book boundary] (-0.05,1.15) rectangle (5.45,2.35);
  \draw[book boundary] (-0.05,-0.25) rectangle (5.45,0.95);
\end{pdfdiagram}
\diagramnote{本图要证明的命题是：推理内存不能按“都是张量”统一处理，模型级、请求级和计划级生命周期必须分开。}
\end{systemdiagram}

底层机器后果也直接来自这张分类。packed weight 需要对齐到 cache line、页或 GPU backend 要求的边界；activation arena 希望少量大块内存重复使用，减少 allocator 和 TLB 抖动；KV cache 希望读流连续，避免 decode attention 每步跨太多页面；workspace pool 要能处理不同算法的最大临时需求，而不是每次 kernel 调用分配释放。

\subsection{activation arena：把生命周期分析落到地址}

activation arena 的目标很明确：预先申请一块或少数几块大内存，让所有短生命周期 activation 在里面复用。它类似传统编译器里的寄存器分配和栈槽复用，只是单位从标量变量变成了大张量。

planner 至少要知道每个中间值的这些字段：

\begin{lstlisting}[numbers=none,caption={arena planner 使用的 value descriptor}]
ValueDesc:
  id
  shape
  dtype
  layout
  byte_size
  alignment
  defined_at
  last_used_at
  materialization_required
  backend_location
\end{lstlisting}

\code{defined_at} 和 \code{last_used_at} 来自 typed graph 或 lowered plan 的拓扑顺序。若一个值被 residual 分支、oracle dump 或多个后续算子消费，它的 \code{last_used_at} 必须延长。若某个 debug 模式要求保存层输出用于 reference 比较，\code{materialization_required} 会阻止 planner 提前覆盖它。若某个值在 GPU device 上，不能被 CPU arena 的地址复用。

一个简化的 arena planner 可以用线性扫描实现：

\begin{lstlisting}[numbers=none,caption={线性扫描式 activation arena 规划}]
events = values sorted by defined_at
free_list = []
active = []

for v in events:
  expire all active values with last_used_at < v.defined_at
  if v is stateful or materialization escapes plan:
    allocate dedicated storage or mark external
    continue
  block = find_free_block(free_list, v.byte_size, v.alignment)
  if no block:
    block = grow_arena(v.byte_size, v.alignment)
  assign v.offset = block.offset
  add v to active
\end{lstlisting}

工业实现会更复杂：要处理 in-place 更新、别名、跨 stream 同步、不同 device、best-fit/first-fit 策略、碎片回收和峰值优化。但最小原则先抓住生命周期：只要 typed graph 告诉我们每个值何时产生、何时最后使用，就可以把多个不重叠生命周期的 activation 放在同一段物理内存上。

\topic{对齐不是语义细节。} CPU kernel 往往希望输入和输出至少按 32 或 64 字节对齐，便于向量 load/store 和 cache line 行为稳定。GPU 后端可能要求更大的 alignment 或特定 pitch。arena planner 若只按 byte size 紧密拼接，最后会把复杂的 unaligned 处理推给每个 kernel。

\topic{清零不是免费操作。} 有些临时 buffer 不需要初始化，因为 kernel 会完整写入；有些 accumulator 或 mask buffer 必须清零。planner 应该把 zero-fill 作为 buffer 属性，而不是让每个算子默认 \code{memset}。否则一个看起来很小的 scratch 可能在长 prompt prefill 时变成可观带宽成本。

\topic{debug/oracle 会改变生命周期。} 打开 oracle 对比时，某些中间结果需要物化并保留到 reference 比较之后。生产模式下能复用的 arena slot，在验证模式下可能不能复用。这不是“debug 模式慢一点”这么简单，而是 planner 必须显式支持不同 materialization policy。

\begin{warningbox}
不要让 kernel 自己偷偷缓存输入指针。arena planner 会复用地址，如果某个异步 kernel 或 debug hook 在生命周期结束后仍持有旧指针，就会读到另一个张量的数据。异步后端必须把 stream/event 同步也纳入生命周期边界。
\end{warningbox}

\subsection{一个小图的 liveness 手算例子}

内存规划不是抽象名词。看一个最小 residual block：

\begin{lstlisting}[numbers=none,caption={最小 residual block 的值流}]
v0 = input_hidden
v1 = rmsnorm(v0)
v2 = linear_qkv(v1)
v3 = attention(v2, kv_cache)
v4 = linear_wo(v3)
v5 = add(v0, v4)
v6 = rmsnorm(v5)
v7 = mlp(v6)
v8 = add(v5, v7)
\end{lstlisting}

若只看最后输出，\code{v1/v2/v3/v4/v6/v7} 都是短生命周期 activation；\code{v0} 必须活到第一次 residual；\code{v5} 必须活到第二次 residual；\code{kv_cache} 是外部状态，不进入 arena。可以把生命周期写成：

\begin{center}
\begin{tabular}{p{0.14\linewidth}p{0.16\linewidth}p{0.18\linewidth}p{0.36\linewidth}}
\toprule
值 & 定义点 & 最后使用 & 规划含义 \\
\midrule
\code{v0} & 输入 & \code{add(v0,v4)} & 不能在 attention 前覆盖 \\
\code{v1} & RMSNorm & \code{linear\_qkv} & QKV 后可复用 \\
\code{v2} & QKV & \code{attention} & attention 后可复用 \\
\code{v3} & attention & \code{linear\_wo} & Wo 后可复用 \\
\code{v5} & residual & \code{add(v5,v7)} & MLP 期间必须保留 \\
\code{v7} & MLP & \code{add(v5,v7)} & 输出后可复用 \\
\bottomrule
\end{tabular}
\end{center}

于是 planner 可以让 \code{v1} 和 \code{v3} 复用同一块 arena，让 \code{v2} 和 \code{v6} 复用另一块，只要它们的 shape/alignment 兼容。这个手算例子也解释了为什么 oracle/debug 会改变峰值内存：若要求保存 \code{v2} 的 Q/K/V 检查点，它的 last use 就不再是 attention，而是 oracle compare。

\begin{lstlisting}[numbers=none,caption={debug policy 对生命周期的影响}]
production:
  v2 last_use = attention

oracle mode:
  v2 last_use = compare_checkpoint_after_layer
  materialization_required = true
\end{lstlisting}

所以运行时不能把“打开 oracle”当成单纯多打印日志。它会改变 materialization policy、arena 峰值和调度同步点。

\subsection{packed weight：把文件布局变成内核布局}

权重是只读的，但这不代表它可以直接以文件布局进入热路径。模型文件常按通用格式存储：张量名、shape、dtype、连续字节。高性能 kernel 需要的却是物理布局：按 tile 重排、按 SIMD 宽度分组、把 scale 和 weight 放在合适距离、补齐尾部、甚至为不同后端生成不同 packed format。

\begin{lstlisting}[numbers=none,caption={从 manifest 到 packed weight 的准备过程}]
for each linear weight tensor:
  verify logical shape and quantization metadata
  choose backend packing format
  allocate aligned packed buffer
  reorder or decode metadata into kernel-friendly panels
  record mapping: logical tensor -> packed descriptor
  keep checksum and source offset for diagnostics
\end{lstlisting}

packed descriptor 不是简单指针。它至少要告诉 kernel：每个 panel 的行列范围、stride、group scale 位置、tail padding、原始 dtype、计算 dtype、对齐和是否允许并行读取。对 int4/group quant 权重来说，scale 与 packed nibble 的关系尤其关键：如果 packing 只重排了 4-bit 数据，却没有同步重排 scale，kernel 会用错量化比例，输出仍然可能看起来“数值正常”。

这一步和传统编译器的后端 lowering 很像。高层 IR 中的 \code{linear(h, Wq)} 是语义；packed weight 是目标机器上的物理表示。语义值相同，不代表物理布局相同。CPU AVX2、AVX-512 和 GPU tensor core 可能需要完全不同的 packed descriptor；runtime 选择 plan 时必须同时选择对应的 packed weight。

\subsection{KV cache planner：跨 token 状态不能当 scratch}

KV cache 是 decoder-only 推理的核心状态。prefill 会批量写入 prompt 中所有 position 的 K/V；decode 每步追加一个 position，并让 attention 读取从开头到当前位置的历史 K/V。它的生命周期跨越整个 session，不能像 activation 一样在 last use 后复用，因为“last use”会随着未来 token 增长而不断后移。

\begin{lstlisting}[numbers=none,caption={KV cache 的逻辑合同}]
append(layer, batch, kv_head, position, k_vector, v_vector)
read(layer, batch, kv_head, range[0..position])

invariants:
  position is monotonically increasing inside one session
  every readable position has been written exactly once
  layout maps logical (layer, head, position, dim) deterministically
  cache capacity is checked before write
\end{lstlisting}

planner 需要先回答容量问题。若直接按最大上下文一次性分配，地址简单，性能稳定，但峰值内存高。若按 page 或 block 增长，短对话更省内存，但 attention kernel 要通过 page table 读取，读流可能变碎。若使用 sliding window 或 KV eviction，必须把语义写进模型合同：并不是所有模型都允许丢弃远处 token；即使允许，position encoding 和 attention mask 也要同步处理。

\begin{center}
\begin{tabular}{p{0.22\linewidth}p{0.28\linewidth}p{0.34\linewidth}}
\toprule
策略 & 优点 & 风险 \\
\midrule
固定最大容量 & 地址简单，kernel 容易优化 & 空间峰值高，短请求浪费大 \\
按 block/page 增长 & 多请求场景更省内存 & page table、碎片和跨页读取增加复杂度 \\
滑动窗口 & 长对话空间可控 & 改变可见上下文，必须匹配模型语义 \\
KV 量化 & 降低 cache 带宽和容量 & attention 精度、scale 存储和反量化成本要验证 \\
\bottomrule
\end{tabular}
\end{center}

KV cache 的物理 layout 要围绕后端访问方式设计。CPU decode 常常让一个线程或一组线程处理部分 head/output 行，连续读取某个 head 的历史 position；此时 \code{position} 和 \code{head_dim} 的局部连续性很重要。GPU kernel 可能按 block/warp 分摊 position 或 head，需要 coalesced load。两者可以使用不同物理 layout，但必须由同一个 logical descriptor 描述，不能让业务层知道后端私有排列。

\begin{systemdiagram}{prefill/decode 与 KV cache 的状态迁移}
\begin{pdfdiagram}[x=1cm,y=1cm]
  \node[book accent node,text width=2.25cm] (tok) at (0,2.65) {tokens\\T prompt};
  \node[book teal node,text width=2.25cm] (prefill) at (3.0,2.65) {prefill plan\\批量算层};
  \node[book purple node,text width=2.25cm] (cache1) at (6.0,2.65) {KV cache\\写 0..T-1};
  \node[book amber node,text width=2.25cm] (logits) at (9.0,2.65) {last logits\\采样入口};

  \node[book navy node,text width=2.25cm] (one) at (0,0.8) {sampled token\\T};
  \node[book teal node,text width=2.25cm] (decode) at (3.0,0.8) {decode plan\\单 token};
  \node[book purple node,text width=2.25cm] (cache2) at (6.0,0.8) {KV cache\\追加 T};
  \node[book amber node,text width=2.25cm] (next) at (9.0,0.8) {next logits\\继续循环};

  \draw[book strong arrow] (tok) -- (prefill);
  \draw[book strong arrow] (prefill) -- (cache1);
  \draw[book strong arrow] (cache1) -- (logits);
  \draw[book strong arrow] (one) -- (decode);
  \draw[book strong arrow] (decode) -- (cache2);
  \draw[book strong arrow] (cache2) -- (next);
  \draw[book weak arrow] (cache1.south) -- (cache2.north);

  \node[book label,text width=9.3cm] at (4.5,-0.55) {last logits 经 sampler 产生下一行的 decode 输入；prefill 写入的历史在 decode 阶段继续可读，decode 每步再追加一个位置。};
\end{pdfdiagram}
\diagramnote{本图要证明的命题是：prefill 和 decode 不是同一形状下的两次执行，而是围绕 KV cache 读写合同形成的两个计划。}
\end{systemdiagram}

\subsection{paged KV allocator：不是 vector push back}

固定最大 KV cache 简单，但并发请求一多就浪费。paged KV cache 把每个 session 的上下文切成固定 token 数的 block，例如每 page 存 \code{16} 或 \code{32} 个 token 的所有层/head K/V。allocator 管理空闲 page，session 的 page table 把逻辑 block 映射到物理 page。

\begin{lstlisting}[numbers=none,caption={paged KV allocator 的状态}]
KvPage:
  page_id
  capacity_tokens
  owner_session
  logical_block
  epoch
  written_mask
  k_bytes
  v_bytes
  scale_bytes_optional

SessionPageTable:
  logical_block -> page_id
  base_position
  window_size_optional
\end{lstlisting}

append 不能只是 \code{push_back}。它要先计算逻辑 block 和 block 内 offset，缺 page 时向 allocator 申请，写入 K/V 和 scale metadata，然后标记 written。读取时必须检查 page 属主、epoch 和 written bit，否则回收后的旧 page 会被当成当前上下文读取。

\begin{lstlisting}[numbers=none,caption={paged KV append/read 的核心检查}]
append(position, k, v):
  block = position / page_tokens
  offset = position % page_tokens
  page = page_table.get_or_allocate(block)
  require page.owner_session == session_id
  require page.epoch == expected_epoch(block)
  store k/v/scale at offset
  page.written_mask.set(offset)

read(position):
  block = position / page_tokens
  offset = position % page_tokens
  page = page_table.lookup(block)
  require page.written_mask.test(offset)
  return K/V view at offset
\end{lstlisting}

这套检查在 debug/oracle 模式必须严格执行；生产模式可以把部分检查移到 page 分配边界，但不能取消语义。否则一旦出现 sliding window 回绕、请求取消、page 复用或 prefix cache 共享，错误会表现成“长上下文偶发漂移”，极难定位。

\subsection{prefill 和 decode：两套计划，两套指标}

prefill 和 decode 的差异不是参数不同，而是计算形态不同。prefill 输入多个 prompt token，Q/K/V projection 和 MLP 更像 GEMM，attention 要处理较大的 token 块；decode 通常一次只处理一个新 token，线性层更像 GEMV 或 small-batch GEMM，attention 主要压力变成读取越来越长的 KV cache。一个 plan 若试图同时覆盖两者，通常会在其中一个阶段浪费机器资源。

\begin{center}
\begin{tabular}{p{0.18\linewidth}p{0.30\linewidth}p{0.30\linewidth}}
\toprule
维度 & prefill & decode \\
\midrule
输入形状 & 多 token prompt & 通常 1 个新 token \\
主要矩阵形态 & GEMM / batched GEMM & GEMV / small-batch GEMM \\
attention 压力 & 块内计算和写 cache & 读历史 KV，随 context 增长 \\
关键指标 & time to first token、prompt 吞吐 & tokens/s、单 token latency、尾延迟 \\
内存行为 & activation 大，arena 峰值高 & activation 小，KV cache 带宽关键 \\
\bottomrule
\end{tabular}
\end{center}

runtime scheduler 需要把请求分成两个阶段，并在每个阶段选择对应 executable plan：

\begin{lstlisting}[numbers=none,caption={prefill/decode 调度主循环}]
plan_prefill = select_plan(stage=prefill, tokens=T, backend)
bind(plan_prefill, session.kv_cache, session.arena)
logits = run(plan_prefill, prompt_tokens)

while not stop:
  next_id = sample(logits, sampler_state)
  session.tokens.append(next_id)
  plan_decode = select_plan(stage=decode, context=session.position, backend)
  bind(plan_decode, session.kv_cache, session.arena)
  logits = run(plan_decode, next_id)
\end{lstlisting}

这里有两个容易忽略的细节。第一，\code{select_plan} 不应该在热路径里重新做 shape inference、fusion 和 packing；它只是在已经编译好的候选计划里选一个。第二，decode plan 可能按 context 长度分档，例如短上下文使用一个简单 kernel，长上下文使用分块 attention 或不同线程划分。分档选择要足够便宜，否则优化收益会被 dispatch 开销吃掉。

\topic{单请求调度和多请求调度不同。} 教材前半部分可以先把一个 session 跑稳；服务端批量推理还要处理多个 session 的 prefill/decode 交错、batching、优先级、取消、超时和 backpressure。无论是否做连续批处理，底层不变量相同：每个 session 的 KV cache 和 token position 不能混，跨请求共享的只有模型级计划和权重。

\topic{线程池不是越满越好。} CPU decode 的单 token latency 常被小 kernel、同步和内存带宽限制。盲目让每个小算子都抢满线程池，可能增加 barrier 和 cache 竞争。scheduler 应该知道哪些 segment 值得并行，哪些 segment 用单线程或少量线程更稳。

\topic{GPU launch 也是成本。} GPU 对大 prefill 很有吸引力，但 decode 的单 token 阶段可能被 kernel launch、host-device 同步和 logits 取回成本影响。runtime 不能只看峰值 FLOPS，必须记录每个阶段的真实事件时间。

\subsection{调度边界：segment 比单算子更适合热路径}

如果 runtime 仍然按单算子逐个 dispatch，就算内存已经规划好，也会留下大量固定开销：查表、虚调用、线程池 barrier、设备同步和 profiler 事件过细。编译器可以把多个已经确定顺序和 buffer 绑定的节点合成 segment。segment 不是融合成一个 kernel，而是一个更粗的调度单元。

\begin{lstlisting}[numbers=none,caption={executable plan 中的 segment 草图}]
Segment layer_0_attention_prefill:
  inputs: hidden, position, kv_cache
  arena_bindings:
    norm_out  -> arena + 0x000000
    q_buffer  -> arena + 0x410000
    k_buffer  -> arena + 0x820000
    v_buffer  -> arena + 0xC30000
  kernels:
    rmsnorm_linear_qkv_fused
    rope_and_kv_append
    attention_prefill
    linear_wo
  sync:
    none on CPU
    stream event before host logits read on GPU
\end{lstlisting}

segment 的价值是把“计划已经决定”这件事落到数据结构里。runtime 执行 segment 时，不再解析 graph，也不重新分配 activation；它只绑定 session 中的外部状态，按预定 kernel 序列推进。profiler 也可以按 segment 汇总，再在需要时打开 kernel 级事件。这样既保留调试能力，又不让热路径永远停留在解释器形态。

这和传统编译器里的基本块、trace 或机器指令调度有相似之处。单条 IR 指令适合分析，最终执行时却希望形成更稳定的机器代码或调度块。AI 推理计划也一样：Graph IR 适合表达语义，segment 适合表达运行期调度。

\subsection{profiler 和 oracle：运行时要留下证据}

没有 profiler，内存 planner 的改动很容易变成主观优化；没有 oracle，融合、量化和 layout 改动很容易悄悄改变 logits。运行时应该把证据作为一等输出，而不是出了问题才临时加日志。

一个最低可用的 profiler 事件可以这样设计：

\begin{lstlisting}[numbers=none,caption={推理 profiler 事件字段}]
event:
  session_id
  stage: load | prefill | decode | sample
  segment_id
  layer
  backend
  start_time_ns
  end_time_ns
  input_tokens
  context_length
  arena_peak_bytes
  workspace_bytes
  kv_cache_bytes
  kernel_name_optional
\end{lstlisting}

这些字段能回答几个关键问题：首 token 慢在 tokenizer、prefill 还是权重 packing；decode 慢在 linear、attention 还是 sampler；上下文变长后 attention 是否按预期变慢；arena 峰值是否符合 memory plan；某个后端切换是否真的改善了当前阶段。

oracle 则回答正确性问题。生产模式不可能每步都跑 reference，但开发和回归测试应该支持可控开关：

\begin{itemize}
  \item 单算子 oracle：比较 RMSNorm、linear、RoPE、attention、MLP 的输出误差。
  \item 层级 oracle：比较某一层输入/输出 hidden state，定位误差扩散起点。
  \item logits oracle：比较 top-k、最大绝对误差、相对误差和采样前分布变化。
  \item 序列 oracle：固定 seed 和 sampler 后，比较生成 token 序列是否满足合同。
  \item 内存 oracle：检查 arena offset 是否重叠错误、KV position 是否越界、debug materialization 是否被提前覆盖。
\end{itemize}

oracle 不是简单打印浮点数。它要知道误差合同：fp32 reference 对 fp16/bf16/int8/int4 后端的容忍度不同；不同 reduction 顺序可能不 bitwise identical；量化 KV cache 会改变 attention 输出，需要单独设误差边界。把这些合同写进测试和报告，才能让优化有可追溯的安全线。

\subsection{硬验收}

读完本节，应该能把 typed graph 到真实请求之间的运行时边界讲清楚：

\begin{itemize}
  \item runtime state 至少包含 session、token buffer、KV cache、activation arena、workspace、sampler 和 profiler/oracle。
  \item 权重、packed weight、activation、workspace、KV cache 和 logits 的生命周期不同，不能放进同一种分配策略。
  \item activation arena 依赖 liveness analysis、last use、alignment、materialization policy 和后端位置。
  \item packed weight 是后端 lowering 的物理产物，必须保存 layout、scale、padding、对齐和诊断映射。
  \item KV cache 是跨 token 状态，容量、layout、page/block、sliding window 和量化策略都必须受模型语义约束。
  \item prefill 和 decode 应该使用不同 executable plan，并用不同指标评估：time to first token、prompt 吞吐、tokens/s 和尾延迟。
  \item segment 把图解释开销移出热路径，profiler 和 oracle 则把性能与正确性证据留在运行时。
\end{itemize}

进入具体执行阶段时，就可以用这份账本判断每个优化是否站得住：它缩短的是哪段生命周期，减少的是哪条地址流，改变的是哪个后端 layout，profiler 应该看到什么，oracle 又必须保证什么。
